A database engine must answer every statement or reject it with an error. Mature
engines check this with differential fuzzers that compare their answers against
a reference engine. Such a fuzzer cannot test a function it never generates, nor
judge one its reference engine lacks. Turso, a SQLite-compatible engine written
in Rust, runs such a fuzzer in continuous integration, yet the single statement
\sql{select LPAD('a', 9223372036854775807, 'b')} terminated the whole process.
We present \tool. It takes the shape of bug reports filed against a different
engine, DuckDB, refills their literals with boundary values, and runs them
against Turso together with a typed SQL generator. Its only oracle needs no
reference engine: one statement must not end the process. \tool found four
previously unknown defects in Turso's string functions, and our two patches that
fix all four were merged upstream. Each defect aborts the engine from one
statement, and the tool regenerates each one from a recorded seed. None of
Turso's own fuzzers can express any of the four, since \sql{lpad},
\sql{rpad} and \sql{repeat} are absent from all of them. At an equal budget of
20{,}000 cases, the typed generator reached three distinct defects and the mined
reports reached two. Running both sources together found the first defect at
case 101. The work taught us that reading an incumbent fuzzer's source is a cheap
and reliable way to decide where to test next.
